\label{sec:abl}
\paragraph{Warmup length at peak LR $10^{-2}$ (H3).} Table~\ref{tab:warm} varies the warmup length. Cosine without warmup reaches \rhval{abl_warmup/warmup+cosine-warmup-0/lr1e-2/val_loss/mean:2} nats, against \rhval{abl_warmup/warmup+cosine-warmup-50/lr1e-2/val_loss/mean:2} with \rhval{const/warm_alt1} warmup steps, \rhval{main/warmup+cosine/lr1e-2/val_loss/mean:2} with the default \rhval{const/warmup_default} (main group) and \rhval{abl_warmup/warmup+cosine-warmup-300/lr1e-2/val_loss/mean:2} with \rhval{const/warm_alt2}. Relative to the no-warmup run, the difference with 50 steps is \rhval{cmp/abl_warmup/warmup+cosine-warmup-50/lr1e-2/val_loss/delta:2} nats (Welch $p=\rhval{cmp/abl_warmup/warmup+cosine-warmup-50/lr1e-2/val_loss/welch_p}$) and with 300 steps \rhval{cmp/abl_warmup/warmup+cosine-warmup-300/lr1e-2/val_loss/delta:2} ($p=\rhval{cmp/abl_warmup/warmup+cosine-warmup-300/lr1e-2/val_loss/welch_p}$). \textbf{H3 is supported.} The registered comparison was against the 100-step default, which lives in the main group and which the comparison tool cannot pool with the ablation group; we therefore tested against the 50- and 300-step runs that flank it and report the 100-step mean from the main table. Beyond roughly 50 steps, more warmup changes little at this rate (the means for 50, 100 and 300 steps are listed above and lie within about two seed standard deviations of each other; we did not test these differences).

\begin{table}[t]\centering\caption{Warmup ablation at peak LR $10^{-2}$ (validation loss).}\label{tab:warm}
\resizebox{\columnwidth}{!}{\begin{tabular}{lcc}
\toprule
Method & val\_loss $\downarrow$ & $n$ \\
\midrule
Warmup+Cosine (warmup=50) & \underline{1.69 $\pm$ 0.01} & 3 \\
Warmup+Cosine (warmup=0) & 2.25 $\pm$ 0.04 & 3 \\
Constant (warmup=100) & 1.81 $\pm$ 0.00 & 3 \\
Constant (warmup=300) & 1.80 $\pm$ 0.00 & 3 \\
Warmup+Cosine (warmup=300) & \textbf{1.67 $\pm$ 0.00} & 3 \\
\bottomrule
\end{tabular}
}\end{table}

\paragraph{Is it the schedule or the warmup?} The same table shows that giving \emph{constant} LR a warmup of 100 or 300 steps lowers its loss at $10^{-2}$ from \rhval{main/constant/lr1e-2/val_loss/mean:2} to \rhval{abl_warmup/constant-warmup-100/lr1e-2/val_loss/mean:2} and \rhval{abl_warmup/constant-warmup-300/lr1e-2/val_loss/mean:2}, so most of the gap between cosine and constant at this rate is attributable to warmup rather than decay; the remaining gap (cosine \rhval{main/warmup+cosine/lr1e-2/val_loss/mean:2} vs.\ constant with warmup 100 at \rhval{abl_warmup/constant-warmup-100/lr1e-2/val_loss/mean:2}) is the decay's contribution. This is a controlled isolation of the warmup factor (same code, one flag), so we state it as a result and not only as a conjecture; why warmup helps (e.g.\ Adam's early second-moment estimate~\cite{liu2019variance}) is \emph{not} tested here.

\begin{table}[t]\centering\caption{Step decay and constant LR with a 100-step warmup at the three main rates (constant at $10^{-2}$ is in Table~\ref{tab:warm}).}\label{tab:fair}
\resizebox{\columnwidth}{!}{\begin{tabular}{llcc}
\toprule
Method & Task & val\_loss $\downarrow$ & $n$ \\
\midrule
Step decay (warmup=100) & lr3e-3 & \textbf{1.71 $\pm$ 0.00} & 3 \\
Constant (warmup=100) & lr3e-3 & \underline{1.74 $\pm$ 0.01} & 3 \\
\midrule
Step decay (warmup=100) & lr1e-3 & \underline{1.83 $\pm$ 0.01} & 3 \\
Constant (warmup=100) & lr1e-3 & \textbf{1.76 $\pm$ 0.02} & 3 \\
\midrule
Step decay (warmup=100) & lr1e-2 & \textbf{1.71 $\pm$ 0.01} & 3 \\
\bottomrule
\end{tabular}
}\end{table}

\paragraph{Fair-warmup baselines and sensitivity.} Table~\ref{tab:fair} adds a 100-step warmup to step decay at all three main rates and to constant LR at $10^{-3}$ and $3\times10^{-3}$ (constant at $10^{-2}$ is in the warmup table). With warmup, step decay's loss at its best rate is \rhval{sens_fair/step-decay-warmup-100/sens/best3/mean:2} (versus \rhval{main/step-decay/sens/best3/mean:2} without) and constant's is \rhval{sens_fair/constant-warmup-100/sens/best3/mean:2} (versus \rhval{main/constant/sens/best3/mean:2}). Their $\text{sens}_3$ drops to \rhval{sens_fair/step-decay-warmup-100/sens/sens3/mean:2} and \rhval{sens_fair/constant-warmup-100/sens/sens3/mean:2}, so \emph{constant LR with warmup has the lowest three-rate sensitivity of any configuration we ran}, slightly below warmup+cosine's \rhval{main/warmup+cosine/sens/sens3/mean:2} (descriptive; no test was run on this comparison, and the ordering could change on a different grid). The cost is a higher best loss than cosine (\rhval{main/warmup+cosine/sens/best3/mean:2}) and Schedule-Free (\rhval{main/schedule-free-adamw/sens/best3/mean:2}). Hence the sensitivity advantage reported for cosine in the previous section belongs largely to warmup, and the schedule shape (decay) mostly buys a lower loss at a given rate. We did not run these fair-warmup variants at $3\times10^{-4}$ or $3\times10^{-2}$, so $\text{sens}_5$ is only available for the four main systems.

\paragraph{Sensitivity to the peak LR (sweep).} The five-rate sweep (Fig.~\ref{fig:lr}, Table~\ref{tab:main}) is itself the second sensitivity study: each system's loss rises on both sides of its best grid point, but the largest rate $3\times10^{-2}$ is still good for warmup+cosine (\rhval{main/warmup+cosine/lr3e-2/val_loss/mean:2}) and Schedule-Free (\rhval{main/schedule-free-adamw/lr3e-2/val_loss/mean:2}), so the optimum of those two lies near or above the upper end of the grid, and a higher rate than $3\times10^{-2}$ was not tried.
